\documentclass[11pt]{article}
\usepackage[T1]{fontenc}
\usepackage{lmodern,microtype}
\usepackage[margin=1in]{geometry}
\usepackage{amsmath,amssymb,amsthm,mathtools}
\usepackage[colorlinks=true,linkcolor=blue,urlcolor=blue]{hyperref}
\setlength{\emergencystretch}{3em}
\newtheorem{theorem}{Theorem}
\newtheorem{lemma}[theorem]{Lemma}
\newcommand{\Q}{\mathbb Q}
\newcommand{\R}{\mathbb R}
\title{Rational arithmetic for Poisson Gaussian noise calibration}
\author{}\date{}
\begin{document}\maketitle

Fix integers $T\ge1$ and $m\ge2$, and rational numbers
$\varepsilon>0$, $0<\delta<1$, $0<q<1$, and $\sigma>0$.
Write
\[
 w_j=\binom mj(1-q)^{m-j}q^j,\qquad
 a_j=\frac{j(j-1)}{2\sigma^2},\qquad E=(m-1)\varepsilon,
 \qquad A=\sum_{j=0}^m w_j e^{a_j}.
\]
The weights and exponents are nonnegative rational numbers. Although $A$
need not be rational, the following sufficient privacy condition can be
established using only rational arithmetic:
\begin{equation}\label{eq:real}
 A^T\le\delta e^E.
\end{equation}
Indeed $A>0$, so taking logarithms gives
$T\log A\le E-\log(1/\delta)$.
The exact moment bound established in the companion note
\emph{Instance-dependent calibration of Poisson-sampled Gaussian noise}
then gives $(\varepsilon,\delta)$-differential privacy.
The mechanism has a fixed public number of optional record slots, with
add/remove adjacency. At each step, independent Bernoulli sampling with
probability $q$ is followed by clipping at radius $C>0$ and independent
Gaussian noise of covariance $C^2\sigma^2 I$. Only the noisy sum is released.
Queries may depend measurably on the full preceding public history. The
bound covers both adjacency directions, an independent public initial seed,
and common randomized postprocessing.

\begin{lemma}[Rational exponential bounds]\label{lem:exp}
For rational $t\ge0$, choose an integer $p\ge0$ with $t\le2^p$ and put
$r=t/2^p\in[0,1]$. For an integer $N\ge1$, define
\[
 P_N(r)=\sum_{k=0}^{N-1}\frac{r^k}{k!},\qquad
 R_N(r)=\frac{r^N}{N!}\frac{N+1}{N},
\]
\[
 L_N(t)=P_N(r)^{2^p},\qquad
 U_N(t)=(P_N(r)+R_N(r))^{2^p}.
\]
These numbers are rational and satisfy
$0<L_N(t)\le e^t\le U_N(t)$.
For each fixed $t,p$, both endpoints converge to $e^t$ as $N\to\infty$.
\end{lemma}
\begin{proof}
The exponential series has nonnegative terms. Moreover,
$(N+\ell)!\ge N!(N+1)^\ell$ for every $\ell\ge0$. Hence
\[
 0\le e^r-P_N(r)
 =\sum_{\ell=0}^{\infty}\frac{r^{N+\ell}}{(N+\ell)!}
 \le\frac{r^N}{N!}\sum_{\ell=0}^{\infty}
       \left(\frac1{N+1}\right)^\ell
 =R_N(r).
\]
Since $P_N(r)\ge1$, raising the inequalities to the integer power $2^p$
preserves their direction. Also $(e^r)^{2^p}=e^t$. Finally,
$R_N(r)\le2/N!\to0$ and $P_N(r)\to e^r$; continuity of a fixed integer
power proves convergence. A sufficient $p$ is found by successively
doubling $1$, using rational comparisons only.
\end{proof}

Apply Lemma~\ref{lem:exp} separately to $a_0,\ldots,a_m,E$, keeping each
chosen range-reduction power fixed as $N$ increases. Define
\[
 B_N=\sum_{j=0}^m w_j U_N(a_j),\qquad D_N=\delta L_N(E).
\]
All these quantities are rational. The proposed test is the exact comparison
\begin{equation}\label{eq:test}
 B_N^T\le D_N.
\end{equation}

\begin{theorem}[Soundness and eventual acceptance]\label{thm:check}
If \eqref{eq:test} holds at any precision $N\ge1$, the mechanism above is
$(\varepsilon,\delta)$-differentially private. If
$A^T<\delta e^E$, then \eqref{eq:test} holds for every sufficiently large $N$.
\end{theorem}
\begin{proof}
Since $w_j\ge0$, termwise upper bounds give $0<A\le B_N$.
Thus acceptance yields
\[
 A^T\le B_N^T\le\delta L_N(E)\le\delta e^E,
\]
which is \eqref{eq:real}. No logarithm or exponential is evaluated by the
test itself.

For termination, finiteness of the sum and Lemma~\ref{lem:exp} give
$B_N\to A$ and $D_N\to\delta e^E$. Consequently
$D_N-B_N^T\to g:=\delta e^E-A^T$.
If $g>0$, choose $N_0$ such that for $N\ge N_0$ the difference between
$D_N-B_N^T$ and $g$ has absolute value less than $g/2$.
Then $D_N-B_N^T>g/2>0$, so the exact comparison succeeds.
\end{proof}

\begin{theorem}[Noise above an admissible threshold]\label{thm:noise}
Put $h=\log(1/\delta)$. Suppose $(m-1)\varepsilon>h$, and let $s_m>0$
be the unique solution of
\[
 T\log A_m(q,1/s_m^2)=(m-1)\varepsilon-h,
 \qquad A_m(q,x)=\sum_{j=0}^m w_j e^{j(j-1)x/2}.
\]
Every rational $\sigma>s_m$ is accepted at all sufficiently large
precisions. Such a rational $\sigma$ exists, and can be chosen in
$(s_m,s_m+\eta)$ for every real $\eta>0$.
\end{theorem}
\begin{proof}
The function $A_m(q,x)$ is continuous and strictly increasing in $x$:
all summands are nondecreasing, and the $j=m$ summand has positive weight
and positive exponent coefficient. It equals $1$ at zero and is unbounded.
Therefore the positive threshold exists uniquely by the intermediate
value theorem. If $\sigma>s_m>0$, then $1/\sigma^2<1/s_m^2$, so
\[
 T\log A_m(q,1/\sigma^2)
 <T\log A_m(q,1/s_m^2)=E-h.
\]
Exponentiating gives $A^T<e^{E-h}=\delta e^E$.
Theorem~\ref{thm:check} applies. Density of the rationals supplies a
rational $\sigma$ in any stated interval.
\end{proof}

For every valid rational request $(T,\varepsilon,\delta,q)$ an admissible
integer $m$ exists, since $(m-1)\varepsilon$ is unbounded as $m$ grows.
Thus rational noise and a finite successful precision always exist.
An exhaustive search that interleaves integer orders, positive rational
noise candidates, and precisions therefore terminates on every valid request.
This does not say that testing precisions forever at an arbitrary fixed
noise terminates: a failing noise, or a noise exactly at a moment threshold,
need not pass. A bounded search may return ``not established.''

The effective factor $c_2=\varepsilon\sigma/
(q\sqrt{T\log(1/\delta)})$ remains a real number even when $\sigma$ is
rational. It can be represented by its defining formula; the rational
privacy test does not require its numerical evaluation. There is no
additional $c_1$ cutoff in this construction. These are sufficient
instance-dependent calibrations, not identifications of the hidden universal
constants in Abadi et al.'s theorem, and not a claim of optimal DP noise.
\end{document}
